% GENERATED FILE. Edit canonical lessons, then run npm run generate:books.
% canonical-source-hash: fnv1a64:0d237f4c9cd3e424
% canonical-lessons: TE-C535-egatali, TE-C535-vekkirincu, TE-C535-mandalincu, TE-C535-niladiyu, TE-C535-mosagincu

\chapter{Teasing and Telling Off}
\label{ch:te-teasing-and-telling-off}
\hlchaptermodality{535}

\begin{chapteropening}
\textbf{By the end of this chapter:} \emph{I can say that someone mocked, scolded, confronted or deceived another.}

Complete the last lesson of chapter 535: I can say that someone mocked, scolded, confronted or deceived another.
\end{chapteropening}

\section[\texorpdfstring{\te{ఎగతాళి} (egatāḷi)}{ఎగతాళి (egatāḷi)}]{\te{ఎగతాళి} (egatāḷi) — mockery; ridicule}

\label{lesson:TE-C535-egatali}



\begin{quote}
Before the new one: say the Telugu for a compromise, then the Telugu for a rift.
\end{quote}

\subsection*{You'll want to know: \te{ఎగతాళి}}
\textbf{\te{ఎగతాళి}} — \emph{egatāḷi} — \textquotedblleft{}mockery; ridicule\textquotedblright{}.

\begin{quote}
\te{కొత్త} \te{అబ్బాయిని} \te{తరగతిలో} \te{ఎగతాళి} \te{చేశారు}.
\end{quote}

\begin{quote}
\emph{kotta abbāyini taragatilō egatāḷi cēśāru.}
\end{quote}

\begin{quote}
They made fun of the new boy in class.
\end{quote}

\te{ఎగతాళి} \te{చేయు}, to make fun of.

\begin{cousinweb}
\textbf{\te{ఎగతాళి}} (\emph{egatāḷi}) is native Dravidian: DEDR 767 lists Telugu \emph{egatāḷi}, \textquotedblleft{}a joke, mockery\textquotedblright{}, and marks Tamil \emph{ekattāḷi} as taken from Telugu.
\end{cousinweb}

\begin{grammarlens}[title={What you've built}]
One more word for describing what you lived and felt.
\end{grammarlens}

\subsection*{Your turn}
\begin{itemize}
\raggedright
  \item \emph{Say it:} \emph{egatāḷi}
  \item \emph{Say it:} \emph{egatāḷi}, once more
  \item \emph{Say it:} \emph{kotta abbāyini taragatilō egatāḷi cēśāru.} — They made fun of the new boy in class.
  \item \emph{Recall:} say the Telugu for a compromise, then the Telugu for a rift, then say \emph{egatāḷi} again
\end{itemize}

\begin{tcolorbox}[breakable,colback=teal!4,colframe=teal!35!black,title={Before you move on}]
What does \te{ఎగతాళి} mean? (\textquotedblleft{}Mockery; ridicule\textquotedblright{}.) Say it once more.
\end{tcolorbox}
\section[\texorpdfstring{\te{వెక్కిరించు} (vekkiriñcu)}{వెక్కిరించు (vekkiriñcu)}]{\te{వెక్కిరించు} (vekkiriñcu) — to mock; to make faces at}

\label{lesson:TE-C535-vekkirincu}



\begin{quote}
Before the new one: say the Telugu for a rift, then the Telugu for mockery.
\end{quote}

\subsection*{You'll want to know: \te{వెక్కిరించు}}
\textbf{\te{వెక్కిరించు}} — \emph{vekkiriñcu} — \textquotedblleft{}to mock; to make faces at\textquotedblright{}.

\begin{quote}
\te{తమ్ముడు} \te{నాలుక} \te{చూపించి} \te{నన్ను} \te{వెక్కిరించాడు}.
\end{quote}

\begin{quote}
\emph{tammuḍu nāluka cūpiñci nannu vekkiriñcāḍu.}
\end{quote}

\begin{quote}
My little brother stuck out his tongue and mocked me.
\end{quote}

Past (I …): \te{వెక్కిరించాను}, \emph{vekkiriñcānu}.

Mocking with faces or by copying someone's voice.

\begin{cousinweb}
\textbf{\te{వెక్కిరించు}} (\emph{vekkiriñcu}) is native Dravidian: DEDR 767 lists Telugu \emph{ekkirincu}, \textquotedblleft{}to make faces at, to mock\textquotedblright{}, in the same article as \te{ఎగతాళి}.
\end{cousinweb}

\begin{grammarlens}[title={What you've built}]
One more word for describing what you lived and felt.
\end{grammarlens}

\subsection*{Your turn}
\begin{itemize}
\raggedright
  \item \emph{Say it:} \emph{vekkiriñcu}
  \item \emph{Say it:} \emph{vekkiriñcu}, once more
  \item \emph{Say it:} \emph{tammuḍu nāluka cūpiñci nannu vekkiriñcāḍu.} — My little brother stuck out his tongue and mocked me.
  \item \emph{Recall:} say the Telugu for a rift, then the Telugu for mockery, then say \emph{vekkiriñcu} again
\end{itemize}

\begin{tcolorbox}[breakable,colback=teal!4,colframe=teal!35!black,title={Before you move on}]
What does \te{వెక్కిరించు} mean? (\textquotedblleft{}To mock; to make faces at\textquotedblright{}.) Say it once more.
\end{tcolorbox}
\section[\texorpdfstring{\te{మందలించు} (mandaliñcu)}{మందలించు (mandaliñcu)}]{\te{మందలించు} (mandaliñcu) — to reprimand; to tell off}

\label{lesson:TE-C535-mandalincu}



\begin{quote}
Before the new one: say the Telugu for mockery, then the Telugu for to mock.
\end{quote}

\subsection*{You'll want to know: \te{మందలించు}}
\textbf{\te{మందలించు}} — \emph{mandaliñcu} — \textquotedblleft{}to reprimand; to tell off\textquotedblright{}.

\begin{quote}
\te{ఆలస్యంగా} \te{వచ్చినందుకు} \te{మేనేజరు} \te{అతన్ని} \te{మందలించారు}.
\end{quote}

\begin{quote}
\emph{ālasyaṅgā vaccinanduku mēnējaru atanni mandaliñcāru.}
\end{quote}

\begin{quote}
The manager told him off for coming late.
\end{quote}

Past (I …): \te{మందలించాను}, \emph{mandaliñcānu}.

Milder than \te{తిట్టు}: a firm word of correction.

\begin{cousinweb}
\textbf{\te{మందలించు}} (\emph{mandaliñcu}) is a Telugu verb for a firm rebuke; its origin is not settled.
\end{cousinweb}

\begin{grammarlens}[title={What you've built}]
One more word for describing what you lived and felt.
\end{grammarlens}

\subsection*{Your turn}
\begin{itemize}
\raggedright
  \item \emph{Say it:} \emph{mandaliñcu}
  \item \emph{Say it:} \emph{mandaliñcu}, once more
  \item \emph{Say it:} \emph{ālasyaṅgā vaccinanduku mēnējaru atanni mandaliñcāru.} — The manager told him off for coming late.
  \item \emph{Recall:} say the Telugu for mockery, then the Telugu for to mock, then say \emph{mandaliñcu} again
\end{itemize}

\begin{tcolorbox}[breakable,colback=teal!4,colframe=teal!35!black,title={Before you move on}]
What does \te{మందలించు} mean? (\textquotedblleft{}To reprimand; to tell off\textquotedblright{}.) Say it once more.
\end{tcolorbox}
\section[\texorpdfstring{\te{నిలదీయు} (niladīyu)}{నిలదీయు (niladīyu)}]{\te{నిలదీయు} (niladīyu) — to call to account; to confront}

\label{lesson:TE-C535-niladiyu}



\begin{quote}
Before the new one: say the Telugu for to mock, then the Telugu for to reprimand.
\end{quote}

\subsection*{You'll want to know: \te{నిలదీయు}}
\textbf{\te{నిలదీయు}} — \emph{niladīyu} — \textquotedblleft{}to call to account; to confront\textquotedblright{}.

\begin{quote}
\te{డబ్బు} \te{గురించి} \te{గ్రామస్థులు} \te{అధికారిని} \te{నిలదీశారు}.
\end{quote}

\begin{quote}
\emph{ḍabbu guriñci grāmasthulu adhikārini niladīśāru.}
\end{quote}

\begin{quote}
The villagers confronted the official about the money.
\end{quote}

Past (I …): \te{నిలదీశాను}, \emph{niladīśānu}.

To stop someone and demand an answer.

\begin{cousinweb}
\textbf{\te{నిలదీయు}} (\emph{niladīyu}) is a native compound: \emph{nilucu}, \textquotedblleft{}to stand, to stop\textquotedblright{} (DEDR 3675), and \emph{tīyu}, \textquotedblleft{}to draw, to pull\textquotedblright{} (DEDR 3407): to pull someone up short.
\end{cousinweb}

\begin{grammarlens}[title={What you've built}]
One more word for describing what you lived and felt.
\end{grammarlens}

\subsection*{Your turn}
\begin{itemize}
\raggedright
  \item \emph{Say it:} \emph{niladīyu}
  \item \emph{Say it:} \emph{niladīyu}, once more
  \item \emph{Say it:} \emph{ḍabbu guriñci grāmasthulu adhikārini niladīśāru.} — The villagers confronted the official about the money.
  \item \emph{Recall:} say the Telugu for to mock, then the Telugu for to reprimand, then say \emph{niladīyu} again
\end{itemize}

\begin{tcolorbox}[breakable,colback=teal!4,colframe=teal!35!black,title={Before you move on}]
What does \te{నిలదీయు} mean? (\textquotedblleft{}To call to account; to confront\textquotedblright{}.) Say it once more.
\end{tcolorbox}
\section[\texorpdfstring{\te{మోసగించు} (mōsagiñcu)}{మోసగించు (mōsagiñcu)}]{\te{మోసగించు} (mōsagiñcu) — to deceive; to cheat}

\label{lesson:TE-C535-mosagincu}



\begin{quote}
Before the new one: say the Telugu for to reprimand, then the Telugu for to call to account.
\end{quote}

\subsection*{You'll want to know: \te{మోసగించు}}
\textbf{\te{మోసగించు}} — \emph{mōsagiñcu} — \textquotedblleft{}to deceive; to cheat\textquotedblright{}.

\begin{quote}
\te{ఆ} \te{దుకాణదారుడు} \te{ముసలివాళ్ళను} \te{మోసగించాడు}.
\end{quote}

\begin{quote}
\emph{ā dukāṇadāruḍu musalivāḷḷanu mōsagiñcāḍu.}
\end{quote}

\begin{quote}
That shopkeeper cheated old people.
\end{quote}

Past (I …): \te{మోసగించాను}, \emph{mōsagiñcānu}.

The noun is \te{మోసం}, deceit.

\begin{cousinweb}
\textbf{\te{మోసగించు}} (\emph{mōsagiñcu}) is built on \te{మోసం}, \textquotedblleft{}deceit\textquotedblright{}, from Sanskrit \emph{moṣa}, \textquotedblleft{}robbery, theft\textquotedblright{} (Monier-Williams), with the Telugu verb ending -\te{గించు}.
\end{cousinweb}

\begin{grammarlens}[title={What you've built}]
One more word for describing what you lived and felt.
\end{grammarlens}

\subsection*{Your turn}
\begin{itemize}
\raggedright
  \item \emph{Say it:} \emph{mōsagiñcu}
  \item \emph{Say it:} \emph{mōsagiñcu}, once more
  \item \emph{Say it:} \emph{ā dukāṇadāruḍu musalivāḷḷanu mōsagiñcāḍu.} — That shopkeeper cheated old people.
  \item \emph{Recall:} say the Telugu for to reprimand, then the Telugu for to call to account, then say \emph{mōsagiñcu} again
\end{itemize}

\begin{tcolorbox}[breakable,colback=teal!4,colframe=teal!35!black,title={Before you move on}]
What does \te{మోసగించు} mean? (\textquotedblleft{}To deceive; to cheat\textquotedblright{}.) Say it once more.
\end{tcolorbox}
